\section*{Bài tập 06 --- Rule-based learning}

\subsection*{Câu 1. PRISM Algorithm}

Mục tiêu phân lớp: \(\text{Job Offer} \in \{\text{Yes}, \text{No}\}\).
Theo PRISM, tại mỗi bước chọn literal có
\[
P(\text{class} \mid \text{literal})=\frac{\#(\text{literal và class})}{\#(\text{literal})}
\]
lớn nhất; nếu bằng nhau thì chọn literal cover nhiều mẫu hơn.

\subsubsection*{a) Sinh luật cho lớp \(\text{Job Offer}=\text{Yes}\)}

\textbf{Rule 1:}
trên tập \(U\) ban đầu (10 mẫu), các literal tốt nhất có accuracy \(=1\):
\[
\text{CGPA}>8:\frac{4}{4},\quad
\text{Practical Knowledge=Very good}:\frac{2}{2},\quad
\text{Communication Skills=Moderate}:\frac{3}{3}.
\]
Chọn literal cover nhiều hơn \(\Rightarrow \text{CGPA}>8\).
\[
R_1:\ \text{IF CGPA}>8\ \text{THEN Job Offer=Yes}.
\]

\textbf{Rule 2:}
loại các mẫu dương đã cover bởi \(R_1\), khi đó literal tốt nhất là
\(\text{Practical Knowledge=Very good}\) với \(\frac{2}{2}=1\):
\[
R_2:\ \text{IF Practical Knowledge=Very good THEN Job Offer=Yes}.
\]

\textbf{Rule 3:}
tiếp tục loại mẫu dương đã cover bởi \(R_2\), còn 1 mẫu dương.
Literal tốt nhất:
\[
\text{Communication Skills=Moderate}:\frac{1}{1}=1.
\]
Suy ra:
\[
R_3:\ \text{IF Communication Skills=Moderate THEN Job Offer=Yes}.
\]

\subsubsection*{b) Sinh luật cho lớp \(\text{Job Offer}=\text{No}\)}

\textbf{Rule 4:}
trên tập ban đầu, literal tốt nhất:
\[
\text{Communication Skills=Poor}:\frac{2}{2}=1.
\]
\[
R_4:\ \text{IF Communication Skills=Poor THEN Job Offer=No}.
\]

\textbf{Rule 5:}
loại mẫu âm đã cover bởi \(R_4\), còn lại literal tốt nhất:
\[
\text{CGPA}<8:\frac{1}{1}=1.
\]
\[
R_5:\ \text{IF CGPA}<8\ \text{THEN Job Offer=No}.
\]

\subsubsection*{c) Rule set cuối cho Câu 1}
\[
\begin{aligned}
R_1&:\ \text{IF CGPA}>8\ \text{THEN Job Offer=Yes},\\
R_2&:\ \text{IF Practical Knowledge=Very good THEN Job Offer=Yes},\\
R_3&:\ \text{IF Communication Skills=Moderate THEN Job Offer=Yes},\\
R_4&:\ \text{IF Communication Skills=Poor THEN Job Offer=No},\\
R_5&:\ \text{IF CGPA}<8\ \text{THEN Job Offer=No}.
\end{aligned}
\]

\subsection*{Câu 2. Xây dựng luật phân lớp bằng PRISM}

Tập dữ liệu có 10 mẫu, nhãn \(\text{Result} \in \{\text{Pass},\text{Fail}\}\).

\subsubsection*{a) Luật cho lớp \(\text{Result}=\text{Pass}\)}

\textbf{Rule 1:}
trên tập ban đầu,
\[
\text{Attendance=Good}:\frac{5}{5}=1
\]
là literal tốt nhất.
\[
R_1:\ \text{IF Attendance=Good THEN Result=Pass}.
\]

\textbf{Rule 2:}
sau khi bỏ các mẫu dương đã được \(R_1\) cover, literal tốt nhất:
\[
\text{Project Submitted=Yes}:\frac{1}{2}=0.5.
\]
Tập con sau literal này còn lẫn lớp, nên tinh chỉnh tiếp:
\[
\text{Attendance=Average}:\frac{1}{1}=1.
\]
Suy ra:
\[
R_2:\ \text{IF Project Submitted=Yes AND Attendance=Average THEN Result=Pass}.
\]

\subsubsection*{b) Luật cho lớp \(\text{Result}=\text{Fail}\)}

\textbf{Rule 3:}
\[
\text{Communication Skill=Poor}:\frac{3}{3}=1
\]
\[
R_3:\ \text{IF Communication Skill=Poor THEN Result=Fail}.
\]

\textbf{Rule 4:}
sau khi loại mẫu âm đã cover bởi \(R_3\), literal tốt nhất:
\[
\text{Attendance=Poor}:\frac{1}{1}=1
\]
\[
R_4:\ \text{IF Attendance=Poor THEN Result=Fail}.
\]

\subsubsection*{c) Viết tập luật IF--THEN cuối cùng}
\[
\begin{aligned}
R_1&:\ \text{IF Attendance=Good THEN Result=Pass},\\
R_2&:\ \text{IF Project Submitted=Yes AND Attendance=Average THEN Result=Pass},\\
R_3&:\ \text{IF Communication Skill=Poor THEN Result=Fail},\\
R_4&:\ \text{IF Attendance=Poor THEN Result=Fail}.
\end{aligned}
\]

\subsubsection*{d) Coverage và Accuracy của từng luật}

Theo đề:
\[
\text{Coverage}=\frac{\#\text{mẫu thỏa antecedent}}{\#\text{tổng mẫu}},\qquad
\text{Accuracy}=\frac{\#\text{mẫu thỏa antecedent và consequent}}{\#\text{mẫu thỏa antecedent}}.
\]

Với tổng mẫu \(=10\), ta có:

\begin{center}
\begin{tabular}{@{}clcc@{}}
\toprule
Luật & Dạng luật & Coverage & Accuracy \\
\midrule
\(R_1\) & IF Attendance=Good THEN Result=Pass & \(5/10=0.50\) & \(5/5=1.00\) \\
\(R_2\) & IF Project Submitted=Yes AND Attendance=Average THEN Result=Pass & \(1/10=0.10\) & \(1/1=1.00\) \\
\(R_3\) & IF Communication Skill=Poor THEN Result=Fail & \(3/10=0.30\) & \(3/3=1.00\) \\
\(R_4\) & IF Attendance=Poor THEN Result=Fail & \(2/10=0.20\) & \(2/2=1.00\) \\
\bottomrule
\end{tabular}
\end{center}

\subsubsection*{e) Conflict resolution}

Trong tập hiện tại, không có mẫu nào cùng lúc khớp các luật có kết luận trái nhau
(một số mẫu có thể khớp nhiều luật nhưng cùng nhãn Fail).
Nếu có conflict, có thể xử lý:
\begin{itemize}
  \item \textbf{Rule ordering:} ưu tiên luật đứng trước trong rule set.
  \item \textbf{Priority list:} gán mức ưu tiên cố định cho từng luật/lớp.
  \item \textbf{Accuracy ordering:} chọn luật có accuracy cao hơn (nếu hòa, xét coverage).
\end{itemize}

\subsubsection*{f) Câu hỏi mở rộng}
\begin{enumerate}
  \item \textbf{PRISM là attribute-oriented hay rule-oriented?}\\
  PRISM là \textbf{rule-oriented}: học trực tiếp từng luật IF--THEN cho từng lớp.
  \item \textbf{PRISM khác Decision Tree ID3 ở điểm nào?}\\
  ID3 tạo cây quyết định phân hoạch toàn bộ dữ liệu theo cấu trúc cây;
  PRISM tạo tập luật độc lập, có thể chồng lấp, và cần cơ chế ưu tiên khi xung đột.
  \item \textbf{Vì sao PRISM thuộc Sequential Covering Algorithm?}\\
  Vì PRISM học \textbf{từng luật một}, sau đó loại các mẫu dương đã được luật đó cover,
  rồi lặp lại cho tới khi không còn mẫu dương chưa cover.
\end{enumerate}

\subsection*{Câu 3. Tính FOIL-Gain và sinh luật bằng FOIL}

Target class: \(\text{Job Offer}=\text{Yes}\).
Tập dữ liệu có \(5\) mẫu:
\[
|P|=2,\quad |N|=3,\quad R_0=\{\}.
\]

FOIL-Gain:
\[
\text{FOIL-Gain}(R,\ell)=p_1\left[\log_2\frac{p_1}{p_1+n_1}-\log_2\frac{p_0}{p_0+n_0}\right].
\]

\subsubsection*{a) Vòng lặp 1: tính gain cho các literal khả dụng}

\begin{center}
\begin{tabular}{@{}lcccc@{}}
\toprule
Literal & \(p_0\) & \(n_0\) & \(p_1,n_1\) & FOIL-Gain \\
\midrule
CGPA \(>9\) & 2 & 3 & \(1,1\) & \(0.3219\) \\
CGPA \(<8\) & 2 & 3 & \(1,1\) & \(0.3219\) \\
Interactiveness = Yes & 2 & 3 & \(2,2\) & \(0.6439\) \\
Practical Knowledge = Good & 2 & 3 & \(2,2\) & \(0.6439\) \\
\bottomrule
\end{tabular}
\end{center}

Các literal có \(p_1=0\) (như Interactiveness=No, Practical Knowledge=Average)
không được chọn.
Literal tốt nhất vòng 1: \(\text{Interactiveness=Yes}\).

\subsubsection*{b) Cập nhật luật sau literal thứ nhất}

\[
R:\ \text{IF Interactiveness=Yes THEN Job Offer=Yes}.
\]
Sau lọc còn:
\[
p_0=2,\quad n_0=2.
\]

\subsubsection*{c) Vòng lặp 2}

Trên tập đã lọc, literal tốt nhất là \(\text{CGPA}<8\) với gain \(=1.0000\).
\[
R_1:\ \text{IF Interactiveness=Yes AND CGPA}<8\ \text{THEN Job Offer=Yes}.
\]
Luật này đã loại hết mẫu âm trong tập hiện tại của \(R_1\), nên kết thúc \(R_1\).

\subsubsection*{d) Lặp để cover các mẫu dương còn lại}

Còn 1 mẫu dương chưa cover, FOIL sinh tiếp:
\[
R_2:\ \text{IF CGPA}>9\ \text{AND Practical Knowledge=Good THEN Job Offer=Yes}.
\]

\subsubsection*{e) Tập luật IF--THEN cuối cùng (FOIL)}
\[
\begin{aligned}
R_1&:\ \text{IF Interactiveness=Yes AND CGPA}<8\ \text{THEN Job Offer=Yes},\\
R_2&:\ \text{IF CGPA}>9\ \text{AND Practical Knowledge=Good THEN Job Offer=Yes}.
\end{aligned}
\]

Khi dự đoán nhị phân, có thể dùng \textbf{default rule}: nếu không khớp luật nào ở trên
thì gán \(\text{Job Offer}=\text{No}\).
\section*{Bài tập 06 --- Rule-based learning}

Dữ liệu PRISM dùng **7 mẫu (ID 4--10)** như tài liệu (trang 2). Mỗi lần PRISM tính
\[
\text{Accuracy}(c \mid a_i)=\frac{\#\text{(thỏa }a_i\text{ và thuộc lớp }c)}{\#\text{(thỏa }a_i)}
\]
trên tập \(U\) hiện tại; chọn cặp có accuracy **cao nhất**, hòa thì chọn cặp **phủ nhiều mẫu hơn**. Sau mỗi luật hoàn chỉnh, **loại khỏi \(U\)** mọi mẫu thỏa luật đó rồi học luật kế tiếp cho cùng lớp.

Với \(T\) toàn tập (\(|T|=7\)):
\[
\text{Coverage}(R)=\frac{\#\text{(antecedent)}}{|T|},\qquad
\text{Accuracy}(R)=\frac{\#\text{(antecedent }\land\text{ consequent)}}{\#\text{(antecedent)}}.
\]

%--------------------------------------------------------------------
\subsection*{PRISM --- lớp Result = Pass}

\textbf{Bước 1} trên \(U=T\): cặp tốt nhất là \textbf{Feature 1 = Good} (accuracy \(1\), phủ \(3\) mẫu, ID 5,\,7,\,10). Luật đã tinh khiết \textbf{chỉ Pass} --- dừng mở rộng luật này.

Loại ba mẫu khỏi \(U\). Trên \(U\) còn lại, \textbf{một} mẫu Pass duy nhất (ID 6).

\textbf{Bước 1} trên \(U\) mới: các cặp có accuracy \(0{,}5\) với độ phủ \(2\) (Feature 1 = Average; Feature 3 = Yes; Feature 4 = Good). Theo quy ước sắp khi hòa, chọn \textbf{Feature 1 = Average}.

Thu hẹp tập xét còn ID 6 và 9; cần bỏ các ca Fail. \textbf{Bước 2}: \textbf{Feature 3 = Yes} cho accuracy \(1\) trên nhánh này (chỉ ID 6).

\medskip
\noindent\textbf{Luật Pass:}
\begin{itemize}\setlength\itemsep{0pt}
  \item \textbf{R1:} IF Feature 1 = Good THEN Result = Pass.
  \item \textbf{R2:} IF Feature 1 = Average AND Feature 3 = Yes THEN Result = Pass.
\end{itemize}

%--------------------------------------------------------------------
\subsection*{PRISM --- lớp Result = Fail}

\textbf{Vòng 1} trên \(U=T\): \textbf{Feature 1 = Poor} và \textbf{Feature 4 = Poor} đều có accuracy \(1\) và cùng phủ **hai** mẫu; theo thứ tự tên thuộc tính (tie-break) chọn \textbf{Feature 1 = Poor} trước.

Loại 4,\,8. Trên \(U=\{5,6,7,9,10\}\), chỉ còn Fail ở ID 9.

\textbf{Vòng 2}: \textbf{Feature 4 = Poor} có accuracy \(1\) (chỉ ID 9 trong \(U\)).

\medskip
\noindent\textbf{Luật Fail:}
\begin{itemize}\setlength\itemsep{0pt}
  \item \textbf{R1:} IF Feature 1 = Poor THEN Result = Fail.
  \item \textbf{R2:} IF Feature 4 = Poor THEN Result = Fail.
\end{itemize}

%--------------------------------------------------------------------
\subsection*{PRISM --- Coverage \& Accuracy (trên $|T|=7$)}

\begin{center}
\begin{tabular}{@{}lcc@{}}\toprule
Luật & Coverage & Accuracy \\\midrule
Pass R1 (Good) & $3/7\approx 0{,}4286$ & $1$ \\
Pass R2 (Average $\land$ Att3=Yes) & $1/7\approx 0{,}1429$ & $1$ \\
Fail R1 (Poor F1) & $2/7\approx 0{,}2857$ & $1$ \\
Fail R2 (Poor F4) & $2/7\approx 0{,}2857$ & $1$ \\\bottomrule
\end{tabular}
\end{center}

(Fail R2 trên toàn \(T\) phủ ID 4 và 9 vì cả hai có Feature 4 = Poor.)

%--------------------------------------------------------------------
\subsection*{(e) Xung đột luật}

Một mẫu có thể thỏa **nhiều** luật cùng lúc (ví dụ ID 4: Fail theo R1 \emph{và} thỏa antecedent của R2 Fail). Cách xử lý thường dùng: \textbf{thứ tự ưu tiên} các luật (list có thứ tự), \textbf{bảng ưu tiên} chọn luật ``đặc biệt'' hơn, hoặc \textbf{sắp theo Accuracy/Coverage} trước khi áp luật đầu tiên khớp.

%--------------------------------------------------------------------
\subsection*{(f) Câu hỏi mở rộng (ý ngắn)}

\textbf{PRISM thiên hướng rule-oriented} (từng luật tối ưu accuracy cục bộ) chứ không xây một cây toàn cục. \textbf{Khác ID3:} ID3 chia đệ quy toàn không gian; PRISM bọc tuần tự (sequential covering) từng lớp, mỗi luật ``gặm'' một phần dữ liệu rồi loại đi. PRISM thuộc họ \textbf{sequential covering} vì luật sau chỉ học trên phần còn lại sau khi các luật trước đã cover.

%--------------------------------------------------------------------
\subsection*{FOIL --- học luật Job Offer = Yes}

\(P=2\) (ID 1,\,2), \(N=3\). Luật rỗng: \(p_0=2,\ n_0=3\).

\[
\mathrm{FOIL\text{-}Gain}=p_1\Bigl[\log_2\frac{p_1}{p_1+n_1}-\log_2\frac{p_0}{p_0+n_0}\Bigr].
\]

\textbf{Vòng 1:} literal có gain lớn nhất là \textbf{Interactiveness = Yes} và \textbf{Practical = Good} (đều \(p_1=2,\ n_1=2\), gain \(\approx 0{,}644\)); chọn \textbf{Interactiveness = Yes} (thứ tự bảng đề ưu tiên trước khi Practical Good).

Sau bước một: \(p_0=n_0=2\) trên tập thỏa \emph{Yes} (ID 1,\,2,\,3,\,5).

\textbf{Vòng 2:} \textbf{CGPA $<8$} có \(p_1=1,\ n_1=0\), gain \(=1{,}0\). Luật 1: \textbf{IF Interactiveness = Yes AND CGPA $<8$ THEN Job Offer = Yes} (chỉ ID 2).

Loại ID 2 khỏi tập dương; còn ID 1. Trên \(\{1,3,4,5\}\): \(p_0=1,\ n_0=3\).

\textbf{Luật 2 vòng 1:} \textbf{CGPA \(\ge 9\)} (gain \(=1{,}0\)). Sau đó \textbf{vòng 2: Practical = Good} thu còn ID 1 thuần dương.

\medskip
\noindent\textbf{Tập luật Yes:}
\begin{itemize}\setlength\itemsep{0pt}
  \item IF Interactiveness = Yes AND CGPA $<8$ THEN Job Offer = Yes.
  \item IF CGPA $\ge 9$ AND Practical Knowledge = Good THEN Job Offer = Yes.
\end{itemize}
